% Source-faithful assembly: Mumford, Abelian Varieties, Appendix II,
% Horizon transcription pages 0272--0281.
\section{Mordell--Weil theorem}
\label{mumford-appendix-II-source}\dcref{M:appII:II}

\begin{theorem}[Mordell--Weil]\label{mumford-thm-mordell-weil}
Let $K$ be a finite extension of $\mathbb Q$ and $X$ an abelian variety over
$K$. Then $X(K)$ is finitely generated.
\source{mumford-abelian-varieties:page-0272}\uses{mumford-prop-weak-mordell-weil,mumford-prop-height-quadratic,mumford-prop-formal-finite-generation}
\end{theorem}
\begin{proposition}[Weak Mordell--Weil]\label{mumford-prop-weak-mordell-weil}
\label{mumford-thm-weak-mordell-weil}
For every $n>1$, $X(K)/nX(K)$ is finite.
\source{mumford-abelian-varieties:page-0272}\uses{mumford-lem-fundamental-mordell-weil,mumford-prop-kummer-deduction}
\end{proposition}
\begin{proposition}[Height pairing]\label{mumford-prop-height-quadratic}
\label{mumford-prop-canonical-height}
There is a symmetric bilinear pairing $\langle\ ,\ \rangle:X(K)^2\to\mathbb R$
with $\langle x,x\rangle\ge0$ and with finitely many $x$ satisfying
$\langle x,x\rangle<C$ for every $C>0$.
\source{mumford-abelian-varieties:page-0272}\uses{mumford-lem-tate-height,mumford-prop-bounded-height-projective}
\end{proposition}
\begin{proposition}[Formal finite-generation criterion]\label{mumford-prop-formal-finite-generation}
\label{mumford-thm-mordell-weil-descent}
If an abelian group $\Gamma$ has finite $\Gamma/n\Gamma$ for every $n>1$ and
has a pairing with the preceding two properties, then $\Gamma$ is finitely generated.
\source{mumford-abelian-varieties:page-0273}
\end{proposition}
\begin{proof}
Choose representatives $x_1,\ldots,x_s$ of $\Gamma/n\Gamma$. For some $C$,
\[
\langle x,x\rangle>C\Rightarrow\langle x-x_i,x-x_i\rangle<2\langle x,x\rangle.\tag{1}
\]
Indeed this follows from
$\langle x-x_i,x-x_i\rangle=\langle x,x\rangle-2\langle x,x_i\rangle+\langle x_i,x_i\rangle$
and Cauchy--Schwarz,
$\langle x,x_i\rangle\le\langle x_i,x_i\rangle^{1/2}\langle x,x\rangle^{1/2}$.
The finite set
$M=\{x_1,\ldots,x_s\}\cup\{x:\langle x,x\rangle<C\}$ generates. Otherwise choose
$x\notin\Gamma_0=\langle M\rangle$ of minimal norm. Then $x-x_i=ny$ for some
$x_i\in M$, and
\[
\langle y,y\rangle=n^{-2}\langle x-x_i,x-x_i\rangle<2n^{-2}\langle x,x\rangle<\langle x,x\rangle.
\]
Thus $y\in\Gamma_0$ and $x=ny+x_i\in\Gamma_0$, contradiction.

\end{proof}

\begin{lemma}[Fundamental lemma]\label{mumford-lem-fundamental-mordell-weil}
For $L=K(n_X^{-1}X(K))$, the composite of the fields obtained by adjoining
all $n$-division points of points of $X(K)$, $L/K$ is finite abelian.
\source{mumford-abelian-varieties:page-0274}
\uses{mumford-cor-division-unramified,mumford-lem-bounded-kummer-extensions}
\end{lemma}
\begin{proposition}[Kummer deduction]\label{mumford-prop-kummer-deduction}
For $G={\rm Gal}(L/K)$,
\[
X(K)/nX(K)\hookrightarrow{\rm Hom}(G,X_n),\qquad
f_x(s)=s(n^{-1}x)-n^{-1}x.
\]
\source{mumford-abelian-varieties:page-0274}\source{mumford-abelian-varieties:page-0275}
\end{proposition}
\begin{proof}
$f_x$ is independent of the division point since $X_n\subset X(K)$;
it is additive, and
$f_x=0\Longleftrightarrow n^{-1}x\in X(K)\Longleftrightarrow x\in nX(K)$.
Thus the map is injective and its target is finite.

\end{proof}

\begin{lemma}[Integral model]\label{mumford-lem-abelian-integral-model}
For a ring of integers $A$ of $K$, after shrinking $Y={\rm Spec}A_s$ there are
a projective $\widetilde X\to Y$, multiplication $\widetilde m$, and section
$\widetilde e$ whose generic fibre is $X$, whose closed fibres are abelian
varieties, and with $\widetilde X(Y)\simeq X(K)$.
\source{mumford-abelian-varieties:page-0275}
\end{lemma}
\begin{proof}[Outline of proof]
Present the variety and its finite diagram of morphisms by finitely many
equations, clear denominators, and shrink the base. Shrink again so all fibres
are smooth and the group diagrams remain valid. Finally take $A_s$ principal;
the projective embedding and
\[
\mathbb P^N_{\widetilde X}(Y)=\{(s_0,\ldots,s_N):s_i\in A_s,\ s_i\text{ coprime}\}/R
\]
show that every generic point extends to a $Y$-point. The omitted smoothness
and projective-point results are the stated standard inputs.
\notready
\end{proof}

\begin{corollary}[Etaleness of division points]\label{mumford-cor-division-etale-integral-model}
If $x\in\widetilde X(Y)$, then $n^{-1}(x)\to\operatorname{Spec}A_s$ is etale
where the residue characteristic does not divide $n$.
\source{mumford-abelian-varieties:page-0276}
\end{corollary}
\begin{proof}
Every component maps onto the base, and the fibre has $n^{2g}$ geometric
points, equal to the degree of the kernel of multiplication by $n$ in the
prime-to-characteristic case. Hence the morphism is etale.

\end{proof}
\begin{corollary}[Unramified division fields]\label{mumford-cor-division-unramified}
There is a finite set $E$ of prime ideals of $A$ such that for $x\in X(K)$ and
$y\in n^{-1}(x)$, $K(y)/K$ is unramified outside $E$.
\source{mumford-abelian-varieties:page-0277}\uses{mumford-cor-division-etale-integral-model}
\end{corollary}

\begin{lemma}[Bounded Kummer extensions]\label{mumford-lem-bounded-kummer-extensions}
The maximal abelian exponent-$n$ extension of $K$ unramified outside fixed $E$
is finite.
\source{mumford-abelian-varieties:page-0277}
\end{lemma}
\begin{proof}
For $F/K$ maximal and $H={\rm Gal}(F/K)$, the sequence
\[
1\to\mu_n\to\overline K^*\xrightarrow{n}\overline K^*\to1
\]
and Hilbert 90 give
\[
K^*/(K^*)^n\simeq{\rm Hom}(H,\mu_n).
\]
Thus the pairing
\[
H\times K^*/(K^*)^n\to\mu_n,\quad
(\sigma,\alpha(K^*)^n)\mapsto\sigma(\sqrt[n]{\alpha})/\sqrt[n]{\alpha}\tag{3}
\]
is nondegenerate, and Kummer theory gives $L=K(\sqrt[n]{B})$ for the
orthogonal subgroup $B$. Locally, unramifiedness outside $E$ is equivalent to
$v(a)\equiv0\pmod n$ for every $v\notin E$. Such elements form finitely many
 classes modulo $K^{*n}$ by finiteness of ideal classes and $E$-units.
\source{mumford-abelian-varieties:page-0278}

\end{proof}

\begin{proof}[Proof of the fundamental lemma]
The division field is abelian of exponent $n$: conjugates of a chosen
$n$-division point differ by elements of $X_n$. By the unramified-division
corollary it is unramified outside a fixed finite set $E$, and the bounded
Kummer extension lemma makes the maximal such extension finite.
\end{proof}

\section{Heights on projective space}\label{mumford-appendix-II-heights-source}\dcref{M:appII:II}
For normalized absolute values,
\[
\log|xy|_v=\log|x|_v+\log|y|_v,\quad
\log|x+y|_v\le\max(\log|x|_v,\log|y|_v)+c_v,
\]
where $c_v=\log2$ archimedean and $0$ otherwise, and
$\sum_vn_v\log|x|_v=0$ for $x\in K^*$.
\begin{lemma}[Projective height]\label{mumford-lem-projective-height}
\[
h(x_0,\ldots,x_n)=\frac1{[K:\mathbb Q]}\sum_vn_v\max_i\log|x_i|_v.\tag{6}
\]
This is well defined on $\mathbb P^n(K)$, invariant under scaling, nonnegative,
and independent of the field containing the coordinates.
\source{mumford-abelian-varieties:page-0279}
\end{lemma}
\begin{proof}
Scaling invariance follows from the product formula. Divide by a nonzero
coordinate to make one coordinate a unit; non-archimedean maxima are then
nonnegative. If $L\supset K$ and $w|v$, then
$K_v\otimes_KL=\bigoplus_{w|v}L_w$ and
$[L:K]n_v=\sum_{w|v}n_w$, which proves field independence.

\end{proof}
\begin{proposition}[Bounded height and degree]\label{mumford-prop-bounded-height-projective}
\label{mumford-lem-northcott-abelian}
For $C,d>0$, finitely many $x\in\mathbb P^n(\overline K)$ satisfy
$h(x)\le C$ and $[K(x):K]\le d$.
\source{mumford-abelian-varieties:page-0280}\source{mumford-abelian-varieties:page-0281}
\end{proposition}
\begin{proof}
Map a point to the coefficients of the norm form
${\rm Nm}(\sum_i x_iT_i)$. The fibres are finite because factorization into
conjugate linear factors is unique. The triangle inequality gives
\[
\max_i\log|x_i'|_v\le d\max_{i,v'}\log|x_i'|_{v'}+d_v,
\]
with $d_v=0$ non-archimedean, hence $h(X_0,\ldots,X_m)\le d'h(x)+d''$.
Bounded-height rational points are finite.

\end{proof}
\begin{proposition}[Equivalent coordinate heights]\label{mumford-prop-equivalent-projective-heights}
Heights $h_1,h_2$ from two coordinate systems satisfy $h_1\sim h_2$.
\source{mumford-abelian-varieties:page-0281}
\end{proposition}
\begin{proof}
For $(x'')=(x')A$, at archimedean places
\[
\max_i\log|x_i''|_v\le\max_i\log|x_i'|_v+\max_{i,j}\log|a_{ij}|_v+\log(n+1),
\]
and omit the last term non-archimedean. Thus $h_2\le h_1+C$; apply the same
argument to $A^{-1}$.

\end{proof}
